% Supplementary information for jsac2_v4 (on-device body digital
% twins). Contents extend the main paper at SI level only --
% no claims live here that aren't summarized in the main text.

\documentclass[10pt,journal,twocolumn]{IEEEtran}

\usepackage{amsmath,amssymb,amsthm,bm}
\usepackage{graphicx}
\usepackage{cite}
\usepackage[T1]{fontenc}
\usepackage{lmodern}
\usepackage{xcolor}
\usepackage{booktabs}
\usepackage{algorithm}
\usepackage{algpseudocode}
\usepackage{hyperref}
\usepackage[capitalise]{cleveref}

\graphicspath{{figures/}{../code/outputs/}}

% Math macros (mirror jsac2_v4.tex).
\newcommand{\rr}{\bm{r}}
\newcommand{\hh}{\bm{h}}
\newcommand{\xx}{\mathbf{x}}
\newcommand{\nhat}{\hat{\bm{n}}}
\newcommand{\khat}{\hat{\bm{k}}}
\newcommand{\etahat}{\hat{\bm{\eta}}}
\newcommand{\Sinc}{S_{\mathrm{inc}}}
\newcommand{\Sab}{S_{\mathrm{ab}}}
\newcommand{\Aab}{A_{\mathrm{ab}}}
\newcommand{\Qabs}{\mathbf{Q}_{\mathrm{ab}}}
\newcommand{\Jmat}{\mathbf{J}}
\newcommand{\Phimat}{\mathbf{\Phi}}
\newcommand{\Fn}{\mathbf{F}_n}
\newcommand{\psin}{\bm{\psi}_n}
\newcommand{\thetab}{\bm{\theta}}
\newcommand{\gammab}{\bm{\gamma}}
\newcommand{\hBS}{\bm{h}_{\mathrm{BS}}}
\newcommand{\hbody}{\bm{h}_{\mathrm{b}}}
\newcommand{\Tepskin}{T_{0}}
\newcommand{\diff}{\mathrm{d}}
\newcommand{\Complex}{\mathbb{C}}
\newcommand{\Reals}{\mathbb{R}}
\newcommand{\ntilde}{\tilde{n}}
\DeclareMathOperator{\RE}{Re}
\DeclareMathOperator{\IM}{Im}

% TAP-style numbering: S1, S2, ... for sections, figures, tables, equations.
\renewcommand{\thesection}{S\arabic{section}}
\renewcommand{\thefigure}{S\arabic{figure}}
\renewcommand{\thetable}{S\arabic{table}}
\renewcommand{\theequation}{S\arabic{equation}}

\theoremstyle{plain}
\newtheorem{proposition}{Proposition}[section]
\theoremstyle{definition}
\newtheorem{definition}[proposition]{Definition}

\title{Supplementary information:\\
On-Device Body Digital Twins for Pose-Aware mmWave Links\\
and Live Exposure Readings}
\author{Robin~Wydaeghe%
  \thanks{R. Wydaeghe is with Ghent University and imec, Ghent, Belgium
  (e-mail: robin.wydaeghe@ugent.be).}}

\begin{document}

\maketitle

% =====================================================================
\section{Approximations 1 and 2: error budgets and proofs}\label{si:approx-proofs}
% =====================================================================

Approximations~1 and~2 of the main paper reduce the coherent absorption
double sum to a single squared norm. Both rest on the same mechanism.
The high refractive index of skin compresses all in-tissue refraction
angles to within $\sim\!10^{\circ}$ of normal. This suppresses the
normal component of the refracted TM field and the angular variation
of the per-path attenuation and phase rates. The proofs follow,
with notation aligned to the main paper.

\subsection{Proof of Approximation~1 (TM polarization direction)}

\begin{proposition}[Approximation~1]\label{prop:S-approx1}
Let $\hat{\bm{t}}_n$ be the unit tangent in the incidence plane, so
$\hat{\bm{k}}_n = \sin\theta_n\,\hat{\bm{t}}_n - \cos\theta_n\,\nhat$.
The exact in-tissue TM direction is $\hat{e}'_{p,n} =
\cos\theta_{t,n}\,\hat{\bm{t}}_n + \sin\theta_{t,n}\,\nhat$ with
$\sin\theta_{t,n} = \sin\theta_n / \ntilde$ from Snell's law in
complex form. Replacing $\hat{e}'_{p,n}$ by the incident-side
$\hat{e}_{p,n} = \cos\theta_n\,\hat{\bm{t}}_n + \sin\theta_n\,\nhat$
in the per-path Fresnel operator $\Fn$ leaves self-terms unchanged
and incurs a non-conservative TM--TM cross-term error of
$O(1/|\ntilde|^2)$.
\end{proposition}

\begin{proof}
Self-terms are exact because $\|\hat{e}_{p,n}\|^2 =
\|\hat{e}'_{p,n}\|^2 = 1$. The TM--TM cross-term contains the inner
product
\begin{equation}\label{eq:eppdot}
  \hat{e}'_{p,n} \cdot \hat{e}'_{p,n'}
  = \cos\theta_{t,n}\cos\theta_{t,n'}
  + \sin\theta_{t,n}\sin\theta_{t,n'}.
\end{equation}
Snell's law gives $\sin\theta_{t,n} = \sin\theta_n / \ntilde =
O(1/|\ntilde|)$ across all incidence angles
$\theta_n \in [0^{\circ}, 90^{\circ}]$. The tangential product
$\cos\theta_{t,n}\cos\theta_{t,n'}$ is therefore the dominant
contribution and the normal-normal product is $O(1/|\ntilde|^2)$.
The same expansion applied to the approximated inner product
$\hat{e}_{p,n} \cdot \hat{e}_{p,n'} = \cos\theta_n\cos\theta_{n'}
+ \sin\theta_n\sin\theta_{n'}$ recovers the same dominant term and
substitutes $\sin\theta_n\sin\theta_{n'}$ (which is $O(1)$) for
$\sin\theta_{t,n}\sin\theta_{t,n'}$ (which is $O(1/|\ntilde|^2)$).
The substitution error on the cross-term is therefore
$O(1/|\ntilde|^2)$.

The physical content is that the in-tissue normal field
$E_\perp \propto \sin\theta_t = O(1/|\ntilde|)$, so its contribution
to $\|\bm{E}^{\mathrm{trans}}\|^2$ is $O(1/|\ntilde|^2)$. For skin
at 28~GHz with $|\ntilde| \approx 4.83$, the bound is
$\lesssim 4\%$. The error shrinks at higher frequencies and on
tissues with larger $|\ntilde|$. Per-tissue values appear in
Table~III of the main paper.
\end{proof}

\subsection{Proof of Approximation~2 (universal depth coupling)}

\begin{proposition}[Approximation~2]\label{prop:S-approx2}
With $\xi_n = \sqrt{\ntilde^2 - 1 + \mu_n^2}$ and the depth
attenuation/phase rates $\alpha_n - j\beta_n = -j k_0 \xi_n$, the
depth coupling
$\Lambda_{nn'} = 1 / [\alpha_n + \alpha_{n'} - j(\beta_{n'} - \beta_n)]$
satisfies
\begin{equation}\label{eq:Gamma-norm}
  \Gamma_{nn'}
    \equiv \frac{\Lambda_{nn'}}{\sqrt{\Lambda_{nn}\Lambda_{n'n'}}}
    = \frac{2\sqrt{\alpha_n \alpha_{n'}}}
           {\alpha_n + \alpha_{n'} - j(\beta_{n'} - \beta_n)},
\end{equation}
with $\Gamma_{nn}=1$. For skin at 28~GHz over all
incidence-angle pairs $(\theta_n, \theta_{n'}) \in
[0^{\circ}, 85^{\circ}]^2$, the substitution $\Gamma_{nn'} \approx 1$
incurs a non-conservative error $|1 - \Gamma_{nn'}| \leq 0.44\%$
with mean $0.17\%$.
\end{proposition}

\begin{proof}
Snell's law in complex form, $\sin\theta_{t,n} = \sin\theta_n/\ntilde$
with $|\ntilde| \approx 4.83$, gives
$|\sin\theta_{t,n}| \leq 1/|\ntilde| \approx 0.21$ across all
incident angles, so $|\theta_{t,n}| \leq \arcsin(0.21) \approx
12^{\circ}$. The tissue refractive index $\ntilde$ therefore
compresses all refracted rays into a $\sim\!10^{\circ}$ cone around
normal. Because $\xi_n = \ntilde \cos\theta_{t,n}$, the variation
of $\xi_n$ across the full incident range is bounded by
$|\xi_n - \ntilde| \leq |\ntilde|(1 - \cos 12^{\circ}) \approx
0.022\,|\ntilde|$. Both $\alpha_n = -k_0 \IM(\xi_n)$ and
$\beta_n = k_0 \RE(\xi_n)$ inherit the same $\sim\!2\%$ relative
spread. The differences $\alpha_n - \alpha_{n'}$ and
$\beta_n - \beta_{n'}$ are therefore second-order in this spread.
Substituting into~\eqref{eq:Gamma-norm} and expanding to first
order in $(\alpha_n - \alpha_{n'})/\bar\alpha$ and
$(\beta_n - \beta_{n'})/\bar\alpha$, with
$\bar\alpha = \tfrac{1}{2}(\alpha_n + \alpha_{n'})$, gives
\begin{equation}
  |1 - \Gamma_{nn'}|
    \leq \frac{1}{2}
    \left(\frac{\alpha_n - \alpha_{n'}}{\bar\alpha}\right)^2
    + \frac{1}{4}
    \left(\frac{\beta_n - \beta_{n'}}{\bar\alpha}\right)^2
    + O(\delta^3),
\end{equation}
with $\delta \sim 0.02$. A numerical sweep across
$(\theta_n, \theta_{n'}) \in [0^{\circ}, 85^{\circ}]^2$ for skin at
28~GHz gives $|1 - \Gamma_{nn'}| \leq 0.44\%$ with mean $0.17\%$.
The bound is non-conservative: it can over- or under-estimate
cross-term magnitudes. The bound is an order of magnitude smaller
than the calibration headroom absorbed by the body-side $\gammab$
regression in the main paper.
\end{proof}

% =====================================================================
\section{Per-path Fresnel transmission operator \texorpdfstring{$\mathbf{F}_n$}{Fn}}\label{si:Fn}
% =====================================================================

The main paper invokes $\Fn$ as rank-2 and polarization-aware
without writing the matrix. The full derivation from Maxwell's
equations (Snell's law via phase matching, TE and TM boundary
conditions on a lossy half-space) is standard~\cite{BornWolf}.

At the surface point $\rr$ on triangle $t$, path $n$ has incidence
cosine $\mu_n = \nhat_t \cdot (-\khat_n)$. For $\mu_n > 0$ the local
TE and TM unit vectors are
\begin{equation}\label{eq:S-TE-TM-basis}
  \hat{e}_{s,n} = \frac{\khat_n \times \nhat_t}
                       {\|\khat_n \times \nhat_t\|},
  \qquad
  \hat{e}_{p,n} = \hat{e}_{s,n} \times \khat_n.
\end{equation}
Continuity of the tangential $\bm{E}$ and $\bm{H}$ fields at $z=0$
gives the Fresnel TE and TM amplitude transmission coefficients
\begin{equation}\label{eq:S-fresnel-t}
  t_{s,n} = \frac{2 \mu_n}{\mu_n + \xi_n},
  \qquad
  t_{p,n} = \frac{2\,\ntilde\,\mu_n}{\ntilde^2\,\mu_n + \xi_n},
\end{equation}
with $\xi_n = \sqrt{\ntilde^2 - 1 + \mu_n^2}$ and $\RE(\xi_n) > 0$.

\begin{definition}[Fresnel transmission operator]\label{def:S-Fn}
The dimensionless rank-2 operator
\begin{equation}\label{eq:S-Fn-def}
  \Fn(\rr) =
  \bigl(t_{s,n}\,\hat{e}_{s,n}\hat{e}_{s,n}^{\!\top}
      + t_{p,n}\,\hat{e}_{p,n}\hat{e}_{p,n}^{\!\top}\bigr)\,H(\mu_n)
\end{equation}
acts on the incident polarization amplitude vector
$\psin = \psi_{s,n}\hat{e}_{s,n} + \psi_{p,n}\hat{e}_{p,n}$ by
projecting onto the local TE and TM directions and scaling each
component by its Fresnel coefficient. The Heaviside factor
$H(\mu_n)$ vanishes on back-facing paths.
\end{definition}

The operator is rank-2 because the incident-side polarization basis
$(\hat{e}_{s,n}, \hat{e}_{p,n})$ spans the 2-D transverse plane to
$\khat_n$. In that basis $\Fn$ is diagonal,
\begin{equation}\label{eq:S-Fn-diag}
  \Fn = \mathrm{diag}(t_{s,n}, t_{p,n}) \cdot H(\mu_n),
\end{equation}
and the diagonal form follows from Approximation~1
(\Cref{prop:S-approx1}). Replacing the refracted TM direction by the
incident TM direction makes the TE and TM channels orthogonal in the
$(\hat{e}_{s,n}, \hat{e}_{p,n})$ basis. The transmitted surface
field is
\begin{equation}\label{eq:S-Etrans}
  \bm{E}_n^{\mathrm{trans}}(\rr, 0)
    = \Fn(\rr)\,\psin\,x_{j(n)}\,e^{-jk_0\khat_n\cdot\rr},
\end{equation}
the per-triangle integrand of the Kirchhoff render in Section~IV of
the main paper and of the exposure operator
$\Qabs(\thetab) = \Jmat^{H}\mathbf{M}(\thetab)\Jmat$.

% =====================================================================
\section{Scene and signal-budget conventions}\label{si:scene}
% =====================================================================

This section lists the scene parameters that the main paper uses
without re-deriving them.

\textbf{Coordinates.} The world frame is right-handed and Z-up. The
user stands at $\bm{r}_u = (30, 0, 1.2)\,$m. The base station sits in
one of three regimes: far-BS at $\bm{r}_\mathrm{BS} = (0, 0, 8)\,$m
(30~m horizontal, 8~m mast), mid-BS at $(20, 0, 5)\,$m, and close-BS
at $(25, 0, 3)\,$m. The phone is held at chest height, $40\,$cm in
front of the body centroid, at world position
$\bm{r}_p = \bm{r}_u + (-0.40, 0, 0.05)\,$m, so its broadside faces
the BS in all three regimes.

\textbf{Body orientation.} The SMPL-X canonical mesh (Y-up, body
facing $+z$) is mapped to the world frame by
$\mathbf{R}_z(-90^{\circ}) \mathbf{R}_x(-90^{\circ})$. The body's
vertical axis then lies along world $+z$, the facing direction along
world $-x$ (toward the BS), the left arm along world $-y$, and the
right arm along world $+y$.

\textbf{Link budget.} The BS transmits at $P_\mathrm{tx} = 43\,$dBm.
The noise floor is $-94\,$dBm ($100\,$MHz bandwidth, $5\,$dB noise
figure). The carrier is $f_c = 28\,$GHz with $100\,$MHz bandwidth,
and the BS $8 \times 8$ URA uses $\lambda/2$ spacing. The main paper's
numerical study (Section~VIII of the main text) runs on the canonical
Sionna RT Munich scene with $16$ traced Rx positions. This SI also
documents a \emph{parametric scalar-loss companion} that the main
paper uses for the comfort-vs.\ rate Pareto and the baseline
comparison: a free-space LOS link with three scalar-loss bands that
enumerate the rest of the propagation budget, namely $35\,$dB
(LOS-attenuated, body is a small perturbation), $55\,$dB
(body-comparable, body-mediated and LOS contributions of equal
magnitude), and $80\,$dB (body-dominant, body channel dominates
direct LOS). The companion is a controlled abstraction. It strips
out the per-scene RT variance and isolates how rate, comfort budget,
and baseline-precoder choice trade off against a single
propagation-budget parameter. Rates are reported as plain Shannon
$R = \log_2(1 + \mathrm{SINR})$ in bps/Hz throughout.

% =====================================================================
\section{VPoser latent-space details}\label{si:vposer}
% =====================================================================

The variational pose prior of \cite{Pavlakos2019VPoser} parameterizes
a learned manifold of plausible body poses. The decoder
$\bm{G}: \Reals^{32} \to \Reals^{63}$ is a 3-layer MLP
($512 \to 512 \to 126$, leaky-ReLU) with a 6-D continuous-rotation
output, decoded to per-joint axis-angle for the 21 SMPL-X body joints
(global orientation and hands excluded). The encoder $\bm{G}^{-1}$
is a parallel 3-layer MLP that returns the mean and log-variance of
an approximate Gaussian posterior. We take the mean for deterministic
AMASS-frame inversion.

\textbf{Pretrained weights.} We load the official V02\_05 checkpoint
(epoch 13), trained on the AMASS corpus~\cite{AMASS2019} with batch
size 128 on the CMU + BMLrub split.

\textbf{Comfort radius.} The latent ball
$\|\bm{z} - \bm{z}_0\| \leq \rho_z$ with $\rho_z = 1.5$ corresponds,
under $p(\bm{z}) = \mathcal{N}(\bm{0}, \bm{I})$, to a prior
log-probability above $\log p(\bm{z}_0) - 1.125$. Empirically, this
threshold rules out the implausible-pose tail while allowing torso
yaws of $\pm 30^{\circ}$ and shoulder rolls of $\pm 25^{\circ}$ on
the seated and standing baselines we evaluate.

\textbf{Encoder round-trip.} An AMASS walk frame projected to its
posterior mean and decoded back has a per-joint reconstruction
error of $\sim\!11^{\circ}$ on average. This matches the
information-bottleneck loss the VAE was trained against, and the
body-side $\gammab$ calibration of the main paper's Section~V
absorbs it entirely.

% =====================================================================
\section{Calibration recovery}\label{si:calibration}
% =====================================================================

\textbf{Tier-D scalar pre-conditioning.} Before the body-side
$K$-mode SVD fit, we fit a single complex scalar $\beta_0$ on the
LOS coefficient and a single real scalar on $\hbody$ to absorb gain
miscalibration. The Tier-B SVD fit then operates on the resulting
Tier-D residual $\bm{r}^{(D)}$.

\textbf{Tier-B operator-SVD basis.} The $K = K_{99}$ leading right
singular vectors of $\Lambda_\mathrm{KH}$ form the basis $\bm{V}_K$.
The empirical operator-spectrum sweep of~\Cref{si:svd} gives
$K_{99} \in \{1,2,4\}$ across the three BS regimes, and the leading
mode absorbs $\geq 93\%$ of the body-channel energy at far-BS.

\textbf{Recovery vs.\ SNR sweep on the plaza geometry.}
\Cref{fig:resid-snr} reports the post-calibration residual
$\rho = \|\hat{\bm{\beta}} - \bm{\gamma}\| / \|\bm{\gamma}\|$ as a
function of UL pilot SNR for path-set sizes
$N \in \{20, 40, 60\}$ on the plaza geometry of the predecessor
manuscript, under both plain LS and ridge LS ($\lambda = 10^{-2}$).
The no-calibration baseline sits at $-4.7\,$dB. Trusting the
in-silico amplitudes leaves a $58\%$ relative error on the per-path
complex amplitudes. Plain LS at $N = 20$ beats the no-cal baseline
above $15\,$dB UL SNR and reaches $-21\,$dB at $30\,$dB UL SNR. At
$N = 60$, the column count approaches the array size ($M = 64$),
and the steering matrix develops small singular values that amplify
pilot noise. Ridge LS then recovers $\sim\!25\,$dB of headroom and
beats the no-cal baseline above $25\,$dB UL SNR. A deployment should
therefore pick $N < M_\mathrm{ant} / 2$, regularize, or merge
correlated paths into clusters.

\begin{figure}[!t]
\centering
\includegraphics[width=\columnwidth]{residual_vs_snr.pdf}
\caption{Post-calibration residual $\rho$ vs.\ UL pilot SNR on the
plaza geometry of the predecessor manuscript, $200$ Monte Carlo
trials per $(\mathrm{SNR}, N)$ cell. Solid: plain LS. Dashed: ridge
LS, $\lambda = 10^{-2}$. Dotted horizontal: no-calibration baseline
at $-4.7\,$dB. The plain-LS failure at $N = 60$ comes from
column-correlated steering matrices, and ridge regularization
recovers $\sim\!25\,$dB.}
\label{fig:resid-snr}
\end{figure}

\textbf{Recovery vs.\ scene loss.} The post-calibration residual
$\rho$ above is invariant under the scene-loss sweep: the Tier-B
fit recovers the body channel equally well across the parametric
scalar-loss bands of~\S\ref{si:scene}. This matches the
operator-spectrum analysis of~\S\ref{si:svd}: the $K_{99}$-mode SVD
basis tracks the dominant channel modes regardless of their per-path
scaling.

% =====================================================================
\section{IMU sensitivity overlay}\label{si:imu}
% =====================================================================

We measure the closed-loop rate loss when the on-device twin
receives a noisy IMU pose estimate. The IMU model adds
i.i.d.\ zero-mean Gaussian noise of standard deviation
$\sigma_\mathrm{joint}$ to each of the $21$ SMPL-X body joint
axis-angle components. The twin renders the body and predicts the
channel at $\thetab + \bm{\epsilon}$, and gradient ascent on the
twin's perceived SINR drives the latent to a $\bm{z}^{\star}$ that
optimizes the perturbed channel, not the true one. The reported
rate loss is the SINR gap between $\bm{z}^{\star}$ and the
noiseless-IMU optimum, evaluated on the user's true channel.

\Cref{fig:imu-overlay} reports the closed-loop SINR loss versus
$\sigma_\mathrm{joint} \in \{0, 2, 4, 8, 16\}^{\circ}$ across the
$18$ Munich scenes, with $5$ Monte-Carlo trials per cell. The
mean rate loss stays under $0.5\,$dB across the swept range. At
the canonical $\sigma_\mathrm{joint} = 4^{\circ}$ operating point
(consumer-grade in-pocket IMU), the mean loss is bounded by
$0.4\,$dB on LOS and $0.1\,$dB on NLOS. At the worst-case
$\sigma_\mathrm{joint} = 16^{\circ}$ (poorly calibrated IMU), the
mean loss is bounded by $0.3\,$dB on both classes. The
$90$th-percentile single-trial loss stays under $\sim\!1.8\,$dB
across the swept range. The closed loop is therefore robust to
IMU drift over the consumer-grade operating range, with no
explicit IMU compensation in the gradient.

The IMU model used here is a steady-state per-tick attitude error,
i.i.d.\ Gaussian on each joint axis-angle component. It captures
the magnitude of the per-tick attitude noise that a consumer-grade
strap-down filter (e.g., Madgwick~\cite{Madgwick2011}) reports at
its operating point, but does not model temporal correlation
(gyro-bias drift) or motion-dependent corrector lag. A follow-up
study tracking these dynamics across a continuous walk would run
a Madgwick or Mahony filter directly in the loop. The per-tick
robustness reported here is a lower bound on what such a fuller
model would report, since the steady-state error envelope
dominates the per-tick gradient distortion.

\begin{figure}[!t]
\centering
\includegraphics[width=\columnwidth]{imu_gamma_overlay.pdf}
\caption{IMU sensitivity sweep. Closed-loop SINR loss
relative to a noiseless-IMU oracle, vs.\ per-joint IMU attitude
error $\sigma_\mathrm{joint}$, mean $\pm 1\sigma$ over $5$
Monte-Carlo trials and the $9 + 9$ LOS/NLOS Munich scenes.}
\label{fig:imu-overlay}
\end{figure}

% =====================================================================
\section{Peak-APD location sensitivity}\label{si:peak-apd}
% =====================================================================

The cohort-side $S_\mathrm{ab}$ reading takes two forms: the
\emph{whole-body integrated} absorbed power $P_\mathrm{abs}
= \int_\Sigma S_\mathrm{ab}\,\mathrm{d}A$, and the \emph{peak local}
$S_\mathrm{ab}$ on a localized body region (the ICNIRP-style
$4\,$cm$^2$ surface-averaged hot spot). The two differ sharply in
their sensitivity to the IMU pose estimate.

\Cref{fig:peak-apd} sweeps both quantities against per-joint
IMU attitude noise $\sigma_\mathrm{joint}$ on the same $18$
Munich scenes. For each scene we evaluate $S_\mathrm{ab}$ on
every visible triangle for the operator's noiseless precoder
(MRT against the cascaded oracle channel), once on the true
pose and once on each of $8$ Monte-Carlo perturbations
$\hat{\bm{\theta}} = \bm{\theta} + \bm{\epsilon}(\sigma)$.

\textbf{Whole-body integrated $P_\mathrm{abs}$ is robust.}
Reporting the perturbed-vs-true ratio
$P_\mathrm{abs}(\hat{\bm{\theta}}) / P_\mathrm{abs}(\bm{\theta})$
as a percentage offset (so $0\%$ is exact recovery, negative
means the reading is biased low):
\begin{itemize}
  \item $\sigma = 2^{\circ}$ (research-grade body suit): mean
    bias $-0.0\%$, $\pm 0.7\%$ standard deviation across
    scenes and trials.
  \item $\sigma = 4^{\circ}$ (Madgwick on a 6-IMU body
    suit~\cite{Madgwick2011}): mean $-0.5\%$, $\pm 1.4\%$.
  \item $\sigma = 8^{\circ}$ (mid-range): mean $-1.2\%$,
    $\pm 2.8\%$.
  \item $\sigma = 16^{\circ}$ (consumer-grade phone-class IMU
    pose, MobilePoser-class accuracy): mean $-3.2\%$,
    $\pm 4.5\%$.
\end{itemize}
The integration over the full body surface averages out the
per-triangle $S_\mathrm{ab}$ redistribution caused by
joint-level pose error: the cohort-side $P_\mathrm{abs}$
reading tracks the true value to within ${\sim}5\%$ across
the entire IMU operating range we consider, with a small
($<4\%$) systematic low-bias at the worst-case noise.

\textbf{Single-triangle peak $S_\mathrm{ab}$ location is fragile.}
The peak triangle's centroid moves by a mean of $52\,$cm at
$\sigma = 2^{\circ}$ and $58\,$cm at $\sigma = 16^{\circ}$
(saturating), with standard deviation across scenes and trials
above $30\,$cm. The mechanism is that several triangles carry
near-equal $S_\mathrm{ab}$ values around the per-tick maximum, so
small pose perturbations flip the argmax across body regions
(forearm$\leftrightarrow$shoulder, hip$\leftrightarrow$knee). The
effect is independent of noise magnitude and therefore not
removable by a better IMU.

\textbf{Implication for cohort framing.} The cohort-monitoring
output supported by the present results is the whole-body
integrated $P_\mathrm{abs}$ (and the time-integrated dose
$\int P_\mathrm{abs}\,dt$ over the user's session), with
anatomical-region resolution explicitly out of scope for the
single-tick reading. Spatial averaging at the regulatory
$4\,$cm$^2$ window may recover a stable peak. We leave that
analysis to a follow-up.

\begin{figure}[!t]
\centering
\includegraphics[width=\columnwidth]{peak_apd_sensitivity.pdf}
\caption{Peak-APD location and whole-body $P_\mathrm{abs}$
sensitivity to IMU pose noise, over $18$ Munich scenes ($9$
LOS, $9$ NLOS), $8$ Monte-Carlo trials per cell, fixed MRT
precoder against the noiseless cascaded channel. (Left)
Euclidean distance between peak-Sab triangle centroids on
perturbed vs.\ true body, mean $\pm 1\sigma$. (Right) Ratio
of whole-body integrated $P_\mathrm{abs}$ on perturbed vs.\
true body, mean $\pm 1\sigma$.}
\label{fig:peak-apd}
\end{figure}

% =====================================================================
\section{Operator spectrum and \texorpdfstring{$K_{99}$}{K99} adaptivity}\label{si:svd}
% =====================================================================

\Cref{fig:rank-cdf} reports the $\epsilon$-effective rank of
$\Qabs$ across $20$ sample positions on the Thelonious phantom for
two channel models: a plaza specular dictionary (LOS plus a ground
bounce plus two facade reflections, $4$ paths per body) and a
3GPP~38.901 UMa-LOS stochastic dictionary ($12$ clusters
$\times$ $5$ subpaths $= 56$ subpaths per body), both at $26\,$GHz
on an $8\times 8$ half-wavelength UPA.

At $\epsilon = 10^{-2}$, the median rank is $4$ for the specular
dictionary and $3$ for the stochastic dictionary, with maxima of
$4$ and $3$ respectively. The specular case drops sharply after
$k = 4$ to the double-precision numerical floor, which follows
from $\mathrm{rank}(\Qabs) \leq B$ with $B = 4$ bounces. The
stochastic case decays smoothly, limited by the
$\sim\!14^{\circ}$ angular resolution of the panel: of the $56$
subpaths, only $3$ to $5$ resolvable BS-side directions exceed
$\epsilon = 10^{-2}$, and a long sub-dominant tail reaches the
numerical floor near $k = 30$.

\Cref{fig:svd-spectrum} reports the singular-value spectrum of
$\Lambda_\mathrm{KH}$ across the three BS regimes of the JSAC2
Munich scene. The far-BS geometry concentrates body-channel
energy in a single mode ($K_{99} = 1$). The close-BS geometry,
where the body subtends a wide solid angle, spreads energy across
$K_{99} = 4$ modes. Both figures cover the same physical regime
(rank-few BS-side excitation of the body) on different array
conventions, so the rank-few claim transfers between them.

\begin{figure}[!t]
\centering
\includegraphics[width=\columnwidth]{rank_cdf.pdf}
\caption{$\epsilon$-effective rank of $\Qabs$ on $20$ Thelonious body
positions at $26\,$GHz with an $8\times 8$ half-wavelength UPA.
Median (blue), per-body traces (grey), $10$--$90$ percentile band
(shaded). Left: plaza specular dictionary, hard cliff at $k = 4$
bounces. Right: 3GPP UMa-LOS stochastic dictionary, smooth decay
limited by array angular resolution. Reproduced from the predecessor
manuscript's rank-check experiment.}
\label{fig:rank-cdf}
\end{figure}

\begin{figure}[!t]
\centering
\includegraphics[width=\columnwidth]{fig2_svd_spectrum.pdf}
\caption{Operator $\Lambda_\mathrm{KH}$ singular-value spectrum
on the JSAC2 Munich scene, normalized per regime. Far-BS is
rank-1. Close-BS is rank-4.}
\label{fig:svd-spectrum}
\end{figure}

% =====================================================================
\section{Strict visibility variant for fingertip features}\label{si:visibility}
% =====================================================================

The default Kirchhoff render of Section~IV of the main paper culls
non-illuminated triangles by a per-path back-face test
$\mu_n = \nhat_t \cdot (-\khat_n) > 0$. The cost is
$O(M \cdot M_\mathrm{tri,vis})$ per pose, with $M_\mathrm{tri,vis}$
the number of front-facing triangles. The back-face test is exact
for convex regions and a faithful approximation for the torso, head,
and limbs at FR2 mesh resolution. It under-counts self-shadowing on
the body's concave features.

\textbf{Strict variant.} For a triangle $t$ to be BS-visible along
path $n$, no other body triangle may block the BS-to-$\bm{c}_t$ ray
segment. The strict visibility flag
$V^{(\mathrm{strict})}_{nt} \in \{0, 1\}$ is the conjunction of the
back-face test with a ray-mesh intersection query:
\begin{equation}\label{eq:V-strict}
  V^{(\mathrm{strict})}_{nt} = H(\mu_n) \cdot
    \mathbf{1}\bigl\{\text{no $t' \neq t$ blocks
      $\bm{c}_t \to \bm{r}_\mathrm{BS}^{(n)}$}\bigr\}.
\end{equation}
The intersection query is run against all $M_\mathrm{tri,total} =
20\,908$ body triangles per (path, candidate-triangle) pair, so the
per-pose cost scales as
$O(M \cdot M_\mathrm{tri,vis} \cdot M_\mathrm{tri,total})$.

\textbf{Cost factor.} The ratio
$M_\mathrm{tri,total} / M_\mathrm{tri,vis} \approx 5$ on the SMPL-X
mesh ($M_\mathrm{tri,total} = 20\,908$ triangles,
$M_\mathrm{tri,vis} \approx 4\,000$ front-facing per BS direction).
The strict variant therefore costs ${\sim}5{\times}$ more per pose
than the back-face default. A bounding-volume hierarchy reduces the
per-query intersection cost to $O(\log M_\mathrm{tri,total})$ but
does not change the asymptotic ratio.

\textbf{When required.} Strict visibility matters when the
self-shadowed feature is comparable in surface area to the
back-face-culled estimate of its visible region. For fingertip,
ear-edge, and inner-thigh features, adjacent body geometry casts a
shadow that the back-face test misses. For torso, head, and
forearm-scale features, the strict and default tests agree to within
the per-triangle sampling resolution. The main paper's numerical
study uses the back-face default. The strict variant is invoked
only for the fingertip-features case.

% =====================================================================
\section{Layered Fabry--P\'{e}rot extension to sub-6~GHz}\label{si:fp}
% =====================================================================

The pseudo-Brewster collapse $\Tepskin \approx 0.54$ at $28\,$GHz
that licenses the single-Fresnel reduction of the main paper does
not hold below ${\sim}10\,$GHz, where layered tissue resonances
(skin, fat, and muscle Fabry--P\'{e}rot dips) introduce oscillations
of $\pm 30\%$ in the per-element transmission coefficient. The
extension recipe replaces the constant $\Tepskin$ by the
flux-averaged transmission
$\bar{T}(f) = 2\int_0^1 T_\mathrm{avg}(\mu, f)\,\mu\,\diff\mu$,
computed once per frequency from a layered-tissue Fresnel sweep.
The body-side $\gammab$ calibration absorbs any residual layered
mismatch at no additional pilot cost.

% =====================================================================
\section{Failure-mode taxonomy}\label{si:failure-modes}
% =====================================================================

The closed loop of Section~VII of the main paper either reaches a
comfort-bounded latent perturbation that achieves the target rate
gain or it does not. Failures fall into four categories, and the
iterate trace alone distinguishes them. \Cref{tab:failure-modes}
lists each category with its diagnostic criterion on the iterates
$\bm{z}_0, \bm{z}_1, \ldots, \bm{z}_K$, the typical underlying
cause, and the mitigation.

\begin{table*}[!t]
\centering
\caption{Failure-mode taxonomy for the comfort-bounded latent
gradient ascent. ``BF'' denotes the brute-force optimum within the
comfort ball $\|\bm{z} - \bm{z}_0\| \leq \rho_z = 1.5$, found by
random sampling. ``Best-iterate'' refers to the iterate
$\bm{z}_{k^\star}$ with the highest measured rate, which the closed
loop returns rather than $\bm{z}_K$.}
\label{tab:failure-modes}
\footnotesize
\begin{tabular}{@{}lp{0.30\textwidth}p{0.27\textwidth}p{0.25\textwidth}@{}}
\toprule
Mode & Diagnostic on iterates & Typical cause & Mitigation \\
\midrule
Stagnant
  & $\|\nabla_{\bm{z}} R\| < \epsilon$ for all $k = 0, \ldots, K$
  & Latent ball is flat at $\bm{z}_0$ (local plateau)
  & Random restart inside comfort ball \\[2pt]
Partial
  & $R(\bm{z}_K) \in [R_\mathrm{BF} - 3, R_\mathrm{BF} - 1]\,$dB
  & FD gradient direction biased by probe noise
  & Larger FD step, more probe averages \\[2pt]
Drift
  & $R(\bm{z}_K) > R(\bm{z}_0)$ but $R(\bm{z}_K) < R(\bm{z}_{k^\star})$
  & Line-search backoff failed to track best-iterate
  & Return best-iterate $\bm{z}_{k^\star}$ \\[2pt]
IMU-floor
  & $R(\bm{z}_K) \approx R(\bm{z}_0)$ at high $\sigma_\mathrm{joint}$
  & FD probe drowns in pose-induced channel-prediction error
  & Require $\sigma_\mathrm{joint} \lesssim 4^{\circ}$ \\
\bottomrule
\end{tabular}
\end{table*}

The first three mitigations are independent: random restarts cure
stagnant ascent, larger FD steps cure partial recovery, and
best-iterate tracking cures drift. The IMU-floor mode is the
operating envelope itself, the bound on $\sigma_\mathrm{joint}$
that the hardware loop must respect. Section~VIII of the main
paper labels each non-reaching trajectory by its dominant mode and
reports the relative frequencies stratified by BS-distance regime.

% =====================================================================
\section{Compliance metrics}\label{si:compliance}
% =====================================================================

\textbf{Whole-body SAR.} The ICNIRP 2020 general-public limit is
$0.08\,$W/kg, averaged over $10\,$g of contiguous tissue and a
$30\,$min time window. At the canonical operating point
($43\,$dBm BS power, $30\,$m range), the whole-body SAR is
$\leq 1.2 \times 10^{-4}\,$W/kg, more than two orders of magnitude
below the limit.

\textbf{Peak local APD.} The ICNIRP 2020 general-public limit
above $6\,$GHz is $4\,$W/m$^{2}$, averaged over $4\,$cm$^{2}$. The
peak-local APD over the body surface is $0.58\%$ of the limit at
the canonical operating point. The local averaging kernel is a
disk of radius $\sqrt{4\,\mathrm{cm}^{2}/\pi} \approx 1.13\,$cm
centred at each surface vertex, and the peak across the body is
reported.

\textbf{Why compliance is not the binding constraint.} At FR2
operating points, the link budget (BS-to-body path loss) limits
$\xx$ before any compliance constraint binds. The exposure reading
of the main paper's Section~VIII is therefore informational: it
reports the per-user dose that the regulatory framework specifies
but does not measure directly.

\bibliographystyle{IEEEtran}
\begin{thebibliography}{9}
\bibitem{BornWolf}
M.~Born and E.~Wolf, \emph{Principles of Optics: Electromagnetic
Theory of Propagation, Interference and Diffraction of Light},
7th~ed. Cambridge, U.K.: Cambridge Univ.\ Press, 1999.

\bibitem{Pavlakos2019VPoser}
N.~Ghorbani, A.~Osman, J.~Romero, D.~Tzionas, and M.~J. Black,
``VPoser: Variational human pose prior for body inverse
kinematics,'' software release, Max Planck Inst.\ Intell.\ Syst.,
T\"ubingen, Germany, 2019. [Online]. Available:
\url{https://github.com/nghorbani/human_body_prior}

\bibitem{AMASS2019}
N.~Mahmood, N.~Ghorbani, N.~F. Troje, G.~Pons-Moll, and M.~J.
Black, ``AMASS: Archive of motion capture as surface shapes,'' in
\emph{Proc.\ IEEE/CVF Int.\ Conf.\ Comput.\ Vis.\ (ICCV)}, Seoul,
Republic of Korea, Oct.~2019, pp.~5442--5451.

\bibitem{Madgwick2011}
S.~O.~H. Madgwick, A.~J.~L. Harrison, and R.~Vaidyanathan,
``Estimation of IMU and MARG orientation using a gradient descent
algorithm,'' in \emph{Proc.\ IEEE Int.\ Conf.\ Rehabil.\ Robot.\
(ICORR)}, Zurich, Switzerland, Jun.--Jul.~2011, pp.~1--7.
\end{thebibliography}

\end{document}
